\documentclass[a4paper,11pt]{article}

% Packages
\usepackage[utf8]{inputenc}
\usepackage[T1]{fontenc}
\usepackage[a4paper,margin=2.5cm]{geometry}
\usepackage{graphicx}            % imagens
\usepackage{booktabs}            % tabelas mais bonitas
\usepackage{enumitem}            % listas
\usepackage{hyperref}            % links clicáveis e índice navegável
\usepackage{fancyhdr}            % cabeçalhos e rodapés
\usepackage{lmodern}
\usepackage{longtable}
\usepackage{float}
\usepackage[export]{adjustbox}
\usepackage{seqsplit}
\usepackage[backend=biber,style=numeric,sorting=none]{biblatex}
%\addbibresource{bib/references.bib}
\usepackage[acronym,nomain,nopostdot]{glossaries}
\makenoidxglossaries

\setcounter{secnumdepth}{0}
\renewcommand{\contentsname}{Índice}
\renewcommand{\figurename}{Figura}
\renewcommand{\tablename}{Tabela}

\hypersetup{
  colorlinks=true,
  linkcolor=black,
  urlcolor=blue,
  pdftitle={Relatorio Cibersegurança},
}

% ---------- Metadados (edita aqui) ----------
\newcommand{\titulodoc}{Relatorio Cibersegurança modulo 3}
\newcommand{\nomeprojeto}{Relatorio Cibersegurança modulo 3}
\newcommand{\autores}{Autores: \\ Raul Casquinha Nº50394 MEET \\ Felisberto Varela Nº50404 MEET \\ Tomás Figueiredo Nº50881 MEIC}
\newcommand{\organizacao}{Instituto Superior de Engenharia de Lisboa}
\newcommand{\versao}{2.0}

% ---------- Cabeçalho e rodapé (não aparecem na capa) ----------
\pagestyle{fancy}
\fancyhf{}
\fancyhead[L]{\nomeprojeto}
\fancyhead[R]{\titulodoc}
\fancyfoot[C]{\thepage}
\renewcommand{\headrulewidth}{0.4pt}

% Espaçamento entre parágrafos (em vez de indentação)
\setlength{\headheight}{14pt}
\setlength{\parindent}{0pt}
\setlength{\parskip}{0.6em}

\begin{document}

% =====================================================
%  CAPA
% =====================================================
\begin{titlepage}
  \thispagestyle{empty}
  \centering

  \vspace*{1cm}
  % Descomenta para colocar uma imagem na capa (ficheiro na mesma pasta):
  % \includegraphics[width=0.35\textwidth]{logo.png}
  % \vspace{2cm}

  {\large \organizacao \par}
  \vspace{3cm}

  {\Huge\bfseries \titulodoc \par}
  \vspace{0.8cm}
  {\LARGE \nomeprojeto \par}

  \vfill

  {\large \autores \par}
  \vspace{1cm}
  {\large \today \par}
  \vspace{1cm}
\end{titlepage}


% =====================================================
%  ÍNDICE
% =====================================================
\pagenumbering{roman}
\tableofcontents
\newpage

% Se quiseres lista de figuras e de tabelas, descomenta:
% \listoffigures
% \listoftables
% \newpage

% Se quiseres glossario com os acronimos, descomenta:
%\glsaddallunused
%\printnoidxglossary[type=\acronymtype,title={Lista de Acrónimos}]
%\newpage

\pagenumbering{arabic}

% ----------------- Relatorio --------------
Boa noite
%\printbibliography[title={Referências}]


% ----------------- Relatorio --------------

\end{document}